\scanpage{173}\mspage{118}
i.e. le foncteur $(*)$
\[
f^*:\mathrm{HOT}_W(S)\longrightarrow\mathrm{HOT}_W(S').
\]

D'autre part, $\mathrm{HOT}_W(S)$ peut être \unsure{décrit|décrite} aussi comme \unsure{localisé|localisée} de $\mathrm{Cat}_{/S}$ \deleted{\unsure{tout}} entier, par rapport au localiseur $W_S^d$, \deleted{qui}\inserted{lequel} sur $\underline{\mathrm{Fib}}\,S$, et $=$ sur $\underline{\mathrm{Liss}}\,S$, coïncide avec $W_S^f$. Mais on se bute ici à la difficulté que le foncteur image inverse
\[
f^*:\mathrm{Cat}_{/S}\longrightarrow\mathrm{Cat}_{/S'},\quad X\longmapsto X\times_S S'\,/S'
\]
n'est \uline{pas} compatible en général (si $S$ n'est \deleted{\unsure{équivalente}}\inserted{isomorphe} à la catégorie ponctuelle, pour $S'$ quelconque) avec les localiseurs $W_S^d,W_{S'}^d$. On le voit p.ex. en prenant $S'=\{s\}$, et en se \unsure{rappélant|rappelant} que $W_S^d$ n'implique \uline{pas}, dans $\mathrm{Cat}_{/S}$ tout entier, la $W$-équivalence par fibres.

Ces ennuis n'ont pas lieu sur $\underline{\mathrm{Liss}}\,S$, où $W_S^d$ et $W_S^f$ coïncident. On peut donc décrire $f^*$ \inserted{sur les $\mathrm{HOT}$} aussi à l'aide de
\[
f^*:\underline{\deleted{\unsure{\mathrm{Fib}}}\mathrm{Liss}}\,S\longrightarrow\underline{\mathrm{Liss}}\,S':\quad X\longrightarrow X\times_S S'\,/S',
\]
en passant aux localisés
\[
f^*:(\underline{\mathrm{Liss}}\,S)(W_S^f)^{-1}\longrightarrow(\underline{\mathrm{Liss}}\,S')(W_{S'}^f)^{-1}.
\]

En d'autres termes :

\uline{Prop 1} Si $X$ dans $\mathrm{Cat}_{/S}$ est \uline{lisse} sur $S$, alors on a un isom. can.
\[
f^*(\mathrm{hot}_{W,S}(X))\simeq\mathrm{hot}_{W,S'}(X'),\quad\text{où }X'=X\times_S S'\,/\unsure{S'}.
\]
% archive sheet number: 172
